% !TEX root = ../main.tex

% 中英标题：\chapter{中文标题}[英文标题]
\chapter{已完成的研究工作及进度安排}

\section{已完成的研究工作}
目前为止，本课题已经完成了蛋白质功能预测研究及多模态预训练与扩散模型进行零样本分类的理论学习，成功获取并整理了相关数据集。按照研究内容所述，上游数据集拟选择UniProt版本2024\_01手动管理和审查条目的部分UniProtKB/Swiss-Prot中不具实验性注释的508245个蛋白质进行大规模数据预训练。下游数据集的构建采用了两种策略：一种是依据时间流划分数据集，另一种是基于序列相似性进行聚类划分。前者确保训练集与测试集中的蛋白质数据无冗余，而后者则确保了不同蛋白质序列间的区分度，从而充分保证了模型在不同条件下的泛用性。这两种策略在处理大规模数据时能够有效提高模型的稳定性和准确性，确保研究结果的可靠性。下表展示了根据不同规则生成的下游数据集的详细统计信息。

\begin{table}[h!]
\centering
\caption{根据不同规则生成的下游数据集统计信息}
\begin{tabular}{@{}lrrrrrr@{}}  % 更改为全部右对齐
\toprule
数据集 & 统计数据 & BPO & CCO & MFO & 全部 \\
\midrule
\multirow{4}{*}{时间流划分数据集} 
 & 训练集 & 47,830 & 42,891 & 31,756 & 57,254 \\
 & 验证集 & 777 & 717 & 684 & 1,432 \\
 & 测试集 & 1,070 & 898 & 408 & 1,435 \\
 & GO标签 & 19,717 & 2,506 & 6,095 & 28,318 \\
\multirow{4}{*}{序列相似性划分数据集} 
 & 训练集 & 44,827 & 40,235 & 29,745 & 54,302 \\
 & 验证集 & 2,378 & 2,091 & 1,510 & 2,823 \\
 & 测试集 & 2,472 & 2,180 & 1,593 & 2,996 \\
 & GO标签 & 19,450 & 2,515 & 5,914 & 27,879 \\
\bottomrule
\end{tabular}
\end{table}

此外，本课题还复现了蛋白质功能预测近年来的经典论文，包括结合多模态学习框架用于蛋白质表示，通过学习蛋白质的序列、结构和功能注释数据来进行多任务学习，其目的在于提高蛋白质相关下游任务的性能，并通过新的空间度量方法量化表示的转移性的MASSA模型\cite{hu2023multimodal}；学习结构域与GO之间的关联，构建功能一致的结构域表示，从而有效地进行蛋白质功能预测的Domain-PFP模型\cite{ibtehaz2023domain}；基于DeepFRI数据集，加入层次图注意力的HEAL模型\cite{gu2023hierarchical}；基于自注意力图池化算法，利用AlphaFold2提供的结构信息来增强模型预测的准确性和通用性的Struct2GO\cite{jiao2023struct2go}模型；使用编码器-解码器架构进行预训练，然后进行微调以学习更有效的蛋白质表示的CFAGO模型\cite{wu2023cfago}；结合预训练的大型语言模型和神经符号模型，利用基因本体（GO）中的公理进行预测的DeepGO-SE模型\cite{kulmanov2024protein}。


\section{进度安排}
本课题的研究时间为2023年9月开始至2025年12月份进行答辩，根据以往项目经验，结合个人的学习和科研攻关能力，为本课题研究计划，制定以下的研究进度安排：

（1）2023年9月至2024年8月：已完成相关理论学习，经典论文复现，获取并整理数据集。

（2）2024年9月至2024年12月：构建双分支变分自编码器（VAE）预训练架构并设计扩散模型，包括向前噪声添加过程、通过逐步去除高斯噪声恢复蛋白质数据的扩散过程，通过预训练学习出适用于下游多标签分类任务的有效潜在空间表示，

（3）2025年1月至2025年3月：设计并实现基于扩散模型的标签类条件分类器，能够处理无标签或稀有标签数据。

（4）2025年3月至2025年5月：进行对比实验，准备中期答辩。

（5）2025年6月至2025年8月：继续优化多标签分类器，确保其在各种复杂条件下的预测结果的稳定性和准确性。

（6）2025年9月至2025年11月：完成论文以及相关材料的撰写，准备答辩。


